\subsection{PERMUTATION GROUPS OF A FINITE SET}

If E is a finite set with n elements, the symmetric group \(S_{E}\) (§ 4, no. 1) is a finite group of order \(n!\) . When E is the interval \([1, n]\) of the set N of natural numbers, the corresponding symmetric group is denoted by \(S_{n}\) ; the symmetric group of any set with n elements is isomorphic to \(S_{n}\) .

DEFINITION 8. Let E be a finite set, \(\zeta \in \mathfrak{S}_{\mathrm{E}}\) a permutation of E and \(\bar{\zeta}\) the subgroup of \(\mathfrak{S}_{\mathrm{E}}\) generated by \(\zeta\) . \(\zeta\) is called a cycle if, under the operation of \(\bar{\zeta}\) on E, there exists one and only one orbit which is not reduced to a single element. This orbit is called the support of \(\zeta\) .

Let \(\zeta\) be a cycle. The support \(\operatorname{supp}(\zeta)\) of \(\zeta\) is the set of \(x \in E\) such that \(\zeta(x) \neq x\) .

The order of a cycle \(\zeta\) is equal to the cardinal of its support. The subgroup \(\bar{\zeta}\) generated by \(\zeta\) operates transitively and faithfully on \(\operatorname{supp}(\zeta)\) . As \(\zeta\) is commutative, \(\operatorname{supp}(\zeta)\) is a principal set under \(\bar{\zeta}\) (no. 6, Example 2) and hence \(\operatorname{Card}(\operatorname{supp}(\zeta)) = \operatorname{Card}(\bar{\zeta})\) .

Lemma 1. Let \((\zeta_{i})_{i\in I}\) be a family of cycles whose supports are pairwise disjoint. Then the \(\zeta_{i}\) are pairwise permutable. Let \(\sigma = \prod_{i\in I}\zeta_{i}\) and \(\overline{\sigma}\) be the subgroup generated by \(\sigma\) . Then \(\sigma(x) = \zeta_{i}(x)\) for \(x\in S_{i}, i\in I\) , and \(\sigma(x) = x\) for \(x\notin \bigcup_{i\in I}S_{i}\) . The mapping \(i\mapsto S_i\) is a bijection of I onto the set of \(\overline{\sigma}\) -orbits not consisting of a single element.

Let \(\zeta\) and \(\zeta'\) be two cycles whose supports are disjoint. If

\[
x \notin \operatorname{supp} (\zeta) \cup \operatorname{supp} \left(\zeta^ {\prime}\right),
\]

then \(\zeta\zeta'(x) = \zeta'\zeta(x) = x\) . If \(x\) belongs to the support of \(\zeta\) , then \(\zeta'(x) = x\) and \(\zeta(x)\) belongs to the support of \(\zeta\) , whence \(\zeta\zeta'(x) = \zeta'\zeta(x) = \zeta(x)\) . Similarly when \(x\) belongs to the support of \(\zeta'\) , then \(\zeta'\zeta(x) = \zeta\zeta'(x) = \zeta'(x)\) . Hence \(\zeta\zeta' = \zeta'\zeta\) . Therefore the \(\zeta_i\) are pairwise permutable and, for \(i \in I\) and \(x \in S_i\) , \(\sigma(x) = \zeta_i(x) \in S_i\) . The mappings \(\sigma\) and \(\zeta_i\) coincide on \(S_i\) , hence \(S_i\) is stable under \(\sigma\) and the subgroup of \(\mathfrak{S}_{S_i}\) generated by the restriction of \(\sigma\) to \(S_i\) operates transitively on \(S_i\) ; therefore \(S_i\) is a \(\overline{\sigma}\) -orbit. As the \(S_i\) are non-empty and are pairwise disjoint, the mapping \(i \mapsto S_i\) is injective. As \(\bigcup_{i} S_i\) is the set of \(x\) such that \(\sigma(x) \neq x\) , every \(\overline{\sigma}\) -orbit not consisting of a single element is one of the \(S_i\) .

PROPOSITION 7. Let E be a finite set and \(\sigma\) a permutation of E. There exists one and only one finite set C of cycles satisfying the two following conditions:

\begin{enumerate}
\def\labelenumi{(\alph{enumi})}
\item
  the supports of the elements of C are pairwise disjoints;
\item
  \(\sigma = \prod_{\zeta \in C} \zeta\) (the elements of \(C\) being pairwise permutable by Lemma 1). Let \(\bar{\sigma}\) be the subgroup generated by \(\sigma\) and let \(S\) be the set of \(\bar{\sigma}\) -orbits not consisting of a single element. For \(s \in S\) , write \(\zeta_s(x) = \sigma(x)\) if \(x \in s\) and \(\zeta_s(x) = x\) if \(x \notin s\) . For all \(s \in S\) , \(\zeta_s\) is a cycle whose support is \(s\) and \(\sigma = \prod_{s \in S} \zeta_s\) , as is seen by applying the two sides to any element of \(E\) . The uniqueness of \(C\) follows from Lemma 1.
\end{enumerate}

DEFINITION 9. A cycle of order 2 is called a transposition.

Let \(x\) and \(y\) be two distinct elements of \(E\) . Let \(\tau_{x,y}\) denote the unique transposition with support \(\{x, y\}\) .

For every permutation \(\sigma\) of E the permutation \(\sigma. \tau_{x,y}. \sigma^{-1}\) is a transposition whose support is \(\{\sigma(x), \sigma(y)\}\) . Hence:

\[
\sigma . \tau_ {x, y}. \sigma^ {- 1} = \tau_ {\sigma (x), \sigma (y)}.\tag{4}
\]

Transpositions thus form a conjugacy class in the group \(\mathfrak{S}_{\mathbb{E}}\) .

PROPOSITION 8. Let E be a finite set. The group \(\mathfrak{S}_{\mathrm{E}}\) is generated by the transpositions.

For every permutation \(\sigma\) let \(F_{\sigma}\) be the set of \(x \in E\) such that \(\sigma(x) = x\) . We show by descending induction on \(p\) that every permutation \(\sigma\) such that \(\operatorname{Card}(F_{\sigma}) = p\) is a product of transpositions. If \(p \geqslant \operatorname{Card}(E)\) , the permutation \(\sigma\) is the identity mapping of \(E\) ; it is the product of the empty family of transpositions. If \(p < \operatorname{Card}(E)\) , suppose that the property is proved for every permutation \(\sigma'\) such that \(\operatorname{Card}(F_{\sigma'}) > p\) . Now \(E - F_{\sigma} \neq \varnothing\) ; let \(x \in E - F_{\sigma}\) and \(y \in \sigma(x)\) . Then \(y \neq x\) and \(y \in E - F_{\sigma}\) . Let \(\sigma' = \tau_{x,y} \cdot \sigma\) . The set \(F_{\sigma'}\) contains \(F_{\sigma}\) and \(x\) and hence \(\operatorname{Card}(F_{\sigma'}) > \operatorname{Card}(F_{\sigma}) = p\) . By the induction hypothesis \(\sigma'\) is a product of transpositions and hence \(\sigma = \tau_{x,y} \cdot \sigma'\) is a product of transpositions.

PROPOSITION 9. Let \(n\) be an integer \(\geqslant 0\) . The group \(\mathfrak{S}_n\) is generated by the transpositions \((\tau_{i,i+1})_{1\leqslant i\leqslant n-1}\) .

By virtue of Proposition 8, it suffices to show that every transposition \(\tau_{p,q}\) , \(1 \leqslant p < q \leqslant n\) , belongs to the subgroup H generated by the \(\tau_{i,i+1}\) , \(1 \leqslant i \leqslant n-1\) . We show this by induction on \(q-p\) . For \(q-p=1\) , it is obvious. If \(q-p>1\) , then (formula (4)) \(\tau_{p,q} = \tau_{q-1,q}\tau_{p,q-1}\tau_{q-1,q}\) . By the induction hypothesis \(\tau_{p,q-1} \in H\) and therefore \(\tau_{p,q} \in H\) .

If \(\sigma \in \mathfrak{S}_n\) , every ordered pair \((i,j)\) of elements of \([1,n]\) such that \(i < j\) and \(\sigma(i) > \sigma(j)\) is called an inversion of \(\sigma\) . Let \(\nu(\sigma)\) denote the number of inversions of \(\sigma\) .

Let P be the additive group of mappings from \(Z^{n}\) to Z. For \(f \in P\) and \(\sigma \in S_{n}\) , let \(\sigma f\) be the element of P defined by

\[
\sigma f (z _ {1}, \dots , z _ {n}) = f (z _ {\sigma (1)}, \dots , z _ {\sigma (n)}).\tag{5}
\]

The action of \(\mathfrak{S}_n\) on \(\mathbf{P}\) thus defined is an operation; for all \(\sigma, \tau \in \mathfrak{S}_n\) and \(f \in \mathbf{P}, ef = f\) and

\[
\begin{array}{r l} \left(\tau (\sigma f)\right) (z _ {1}, \dots , z _ {n}) & = \sigma f (z _ {\tau (1)}, \dots , z _ {\tau (n)}) = f (z _ {\tau \sigma (1)}, \dots , z _ {\tau \sigma (n)}) \\ & = ((\tau \sigma) f) (z _ {1}, \dots , z _ {n}). \end{array}
\]

Formula (5) shows that \(\sigma(-f) = -\sigma f\) for \(\sigma \in \mathfrak{S}_n\) and \(f \in P\) .

Let \(p\) be the element of P defined by

\[
p (z _ {1}, \dots , z _ {n}) = \prod_ {i <   j} (z _ {j} - z _ {i}).\tag{6}
\]

Lemma 2. \(p \neq 0\) and \(\sigma p = (-1)^{\nu(\sigma)}\) for \(\sigma \in \mathfrak{S}_n\) .

\(p(1,2,\ldots,n) = \prod_{i < j}(j - i)\neq 0\) and hence \(p\neq 0\) . On the other hand, if \(\sigma \in \mathfrak{S}_n\) , then

\[
\sigma p (z _ {1}, \dots , z _ {n}) = p (z _ {\sigma (1)}, \dots , z _ {\sigma (n)}) = \prod_ {i <   j} (z _ {\sigma (j)} - z _ {\sigma (i)}).
\]

Let C be the set of ordered pairs \((i,j)\) such that \(1 \leqslant i \leqslant n\) , \(1 \leqslant j \leqslant n\) , i \textless{} j. A permutation \(\theta\) is defined on C by setting \(\theta(i,j) = (\sigma(i), \sigma(j))\) if \((i,j)\) is not an inversion, \(\theta(i,j) = (\sigma(j), \sigma(i))\) if \((i,j)\) is an inversion. This implies \(\sigma p = (-1)^{\nu(\sigma)} p\) .

THEOREM 2. Let E be a finite set. There exists one and only one homomorphism \(\varepsilon\) from \(\mathfrak{S}_n\) to the multiplicative group \(\{-1, +1\}\) such that \(\varepsilon(\tau) = -1\) for every transposition \(\tau\) .

The uniqueness follows from Proposition 8. We show the existence. By transporting the structure, it may be assumed that \(\mathbf{E} = [1, n]\) . Using the above notation, let \(\varepsilon(\sigma) = (-1)^{\nu(\sigma)}\) . Then (Lemma 2)

\[
\sigma (\sigma^ {\prime} p) = \sigma (\varepsilon (\sigma^ {\prime}) p) = \varepsilon (\sigma^ {\prime}) (\sigma p) = \varepsilon (\sigma^ {\prime}) \varepsilon (\sigma) p.
\]

On the other hand,

\[
\sigma (\sigma^ {\prime} p) = (\sigma \sigma^ {\prime}) p = \varepsilon (\sigma \sigma^ {\prime}) p.
\]

As \(p \neq -p\) , it follows that \(\varepsilon(\sigma\sigma') = \varepsilon(\sigma)\varepsilon(\sigma')\) and thus \(\varepsilon\) is a homomorphism. We now show that, for every transposition \(\tau, \varepsilon(\tau) = -1\) . \(\nu(\tau_{n-1,n}) = 1\) , whence \(\varepsilon(\tau_{n-1,n}) = -1\) . As every transposition \(\tau\) is conjugate to \(\tau_{n-1,n}\) and the group \(\{-1, +1\}\) is commutative, \(\varepsilon(\tau) = \varepsilon(\tau_{n-1,n}) = -1\) .

DEFINITION 10. In the notation of Theorem 2, the number \(\varepsilon(\sigma)\) (also denoted \(\varepsilon_{\sigma}\) ) is called the signature of the permutation \(\sigma\) . The kernel of the homomorphism \(\varepsilon\) is called the alternating group of E.

\(\sigma\) is called even (resp. odd) if \(\varepsilon(\sigma) = 1\) (resp. \(\varepsilon(\sigma) = -1\) ). The alternating group of E is denoted by \(\mathfrak{A}_{\mathbb{E}}\) . It is a normal subgroup of \(\mathfrak{S}_{\mathbb{E}}\) . When \(\mathbf{E} = [1, n]\) , it is simply denoted by \(\mathfrak{A}_{n}\) . When the cardinal \(n\) of E is \(\geqslant 2\) , it is a subgroup of index 2 and hence of order \(n! / 2\) . It can be shown that, for \(n = 3\) or \(n \geqslant 5\) , the group \(\mathfrak{A}_{n}\) is a simple group (cf.~Exercise 16).

Example. If \(\sigma\) is a cycle of order \(d\) , then

\[
\varepsilon (\sigma) = (- 1) ^ {d - 1}.
\]

The number of inversions of the permutation

\[
(1, 2, 3, \dots , d) \mapsto (d, 1, 2, \dots , d - 1)
\]

is equal to \(d - 1\) .

